\paragraph{fold-\texorpdfstring{$\pi$}{pi} compact 宿主的 finite-rank 耗尽塔与 finite-depth exactness obstruction}
沿用定理 \ref{thm:xi-time-part9zbc-foldpi-oddlayer-positive-moment-formula}、
\ref{thm:xi-time-part9zbc-foldpi-hankel-strict-positivity}、
\ref{thm:xi-time-part9zbe-foldpi-canonical-compact-jacobi}
与 \ref{thm:xi-time-part9zbe-foldpi-compact-determinant-sampling} 的记号。
记
\[
\alpha_r:=\mathfrak a_r^{(\pi)}
=
\frac{B_{2r}}{r(2r-1)}
\bigl(\varphi^{-1}+\varphi^{2r-3}\bigr),
\qquad
u_k:=(-1)^k\frac{\alpha_{k+1}}{(2k)!}
\qquad (r\ge 1,\ k\ge 0).
\]
由前述结果，\((u_k)_{k\ge 0}\) 正是 compact Jacobi 宿主 \(J_\varphi\) 的矩列，因而 H.111 的 odd Bernoulli tail 不仅对应一个无限秩谱宿主，而且还可被一列 finite-rank scalar hosts 从下方单调耗尽。关键点在于：这一耗尽塔对每一层都严格逼近，但在任何有限深度都不会稳定。

\begin{theorem}[H.111 奇层矩的 finite-rank 耗尽塔]\label{thm:xi-time-part9zbg-foldpi-finite-rank-moment-exhaustion}
对每个 \(N\ge 1\)，定义截断测度
\[
\mu_{\varphi,N}
:=
\sum_{n=1}^{N}\frac{\varphi^{-1}}{\pi^2n^2}
\left(
\delta_{(4\pi^2n^2)^{-1}}
+
\delta_{\varphi^2(4\pi^2n^2)^{-1}}
\right),
\]
以及其矩
\[
u_k^{(N)}:=\int y^k\,\dd\mu_{\varphi,N}(y)
\qquad (k\ge 0).
\]
则对任意 \(k\ge 0\)，有显式闭式
\[
u_k^{(N)}
=
4(2\pi)^{-2k-2}
\bigl(\varphi^{-1}+\varphi^{2k-1}\bigr)
\sum_{n=1}^{N}n^{-2k-2},
\]
并且
\[
0<u_k^{(1)}<u_k^{(2)}<\cdots<u_k^{(N)}\uparrow u_k.
\]
更精确地，
\[
u_k-u_k^{(N)}
=
4(2\pi)^{-2k-2}
\bigl(\varphi^{-1}+\varphi^{2k-1}\bigr)
\sum_{n>N}n^{-2k-2}.
\]
\leanverified{paper\_xi\_time\_part9zbg\_foldpi\_finite\_rank\_moment\_exhaustion}
\end{theorem}
\begin{proof}
由 \(\mu_{\varphi,N}\) 的定义，
\[
u_k^{(N)}
=
\sum_{n=1}^{N}\frac{\varphi^{-1}}{\pi^2n^2}
\left(
\frac{1}{(4\pi^2n^2)^k}
+
\frac{\varphi^{2k}}{(4\pi^2n^2)^k}
\right).
\]
整理即得
\[
u_k^{(N)}
=
4(2\pi)^{-2k-2}
\bigl(\varphi^{-1}+\varphi^{2k-1}\bigr)
\sum_{n=1}^{N}n^{-2k-2}.
\]
又因为
\[
\mu_{\varphi,N+1}-\mu_{\varphi,N}
=
\frac{\varphi^{-1}}{\pi^2(N+1)^2}
\left(
\delta_{(4\pi^2(N+1)^2)^{-1}}
+
\delta_{\varphi^2(4\pi^2(N+1)^2)^{-1}}
\right)
\]
是非零正测度，故对每个 \(k\ge 0\) 都有
\[
u_k^{(N+1)}-u_k^{(N)}
=
\int y^k\,\dd(\mu_{\varphi,N+1}-\mu_{\varphi,N})(y)>0.
\]
于是 \((u_k^{(N)})_{N\ge 1}\) 严格递增。
其上界 \(u_k\) 则由定理 \ref{thm:xi-time-part9zbc-foldpi-oddlayer-positive-moment-formula} 给出，而极限与尾项公式都由
\[
\mu_\varphi-\mu_{\varphi,N}
=
\sum_{n>N}\frac{\varphi^{-1}}{\pi^2n^2}
\left(
\delta_{(4\pi^2n^2)^{-1}}
+
\delta_{\varphi^2(4\pi^2n^2)^{-1}}
\right)
\]
直接得到。证毕。
\end{proof}

\begin{theorem}[finite-depth scalar host 的 Jacobi--Fredholm 实现]\label{thm:xi-time-part9zbg-foldpi-finite-rank-jacobi-fredholm-realization}
对每个 \(N\ge 1\)，存在唯一的 \(2N\times 2N\) Jacobi 矩阵
\[
J_{\varphi,N}
=
\begin{pmatrix}
\beta_0^{(N)} & \alpha_1^{(N)} \\
\alpha_1^{(N)} & \beta_1^{(N)} & \alpha_2^{(N)} \\
& \ddots & \ddots & \ddots \\
&& \alpha_{2N-1}^{(N)} & \beta_{2N-1}^{(N)}
\end{pmatrix},
\qquad
\alpha_j^{(N)}>0,\ \beta_j^{(N)}\in\RR,
\]
以及循环向量 \(v_0^{(N)}\)，使得
\[
u_k^{(N)}=\langle v_0^{(N)},(J_{\varphi,N})^k v_0^{(N)}\rangle
\qquad (k\ge 0).
\]
其谱精确为
\[
\sigma(J_{\varphi,N})
=
\left\{\frac{1}{4\pi^2n^2}:1\le n\le N\right\}
\cup
\left\{\frac{\varphi^2}{4\pi^2n^2}:1\le n\le N\right\}.
\]
相应的 cyclic resolvent 为
\[
m_{\varphi,N}(z)
:=
\langle v_0^{(N)},(I-zJ_{\varphi,N})^{-1}v_0^{(N)}\rangle
=
\int\frac{\dd\mu_{\varphi,N}(y)}{1-yz},
\]
即
\[
m_{\varphi,N}(z)
=
\sum_{n=1}^{N}\frac{\varphi^{-1}}{\pi^2n^2}
\left(
\frac{1}{1-z/(4\pi^2n^2)}
+
\frac{1}{1-\varphi^2z/(4\pi^2n^2)}
\right).
\]
因此 \(m_{\varphi,N}\) 是恰有 \(2N\) 个极点的有理 Stieltjes 函数，其 canonical scalar continued fraction 属于 terminating Stieltjes class；关于有限点支撑正交测度与对应 Jacobi/continued-fraction 实现的标准背景，可参见 \href{https://dlmf.nist.gov/18.2}{DLMF 18.2}。
\end{theorem}
\begin{proof}
在有限维 Hilbert 空间 \(L^2(\mu_{\varphi,N})\) 上考虑乘法算子
\[
(M_yf)(y):=y f(y).
\]
由于 \(\mu_{\varphi,N}\) 支撑于 \(2N\) 个互异点，\(L^2(\mu_{\varphi,N})\) 的维数即为 \(2N\)，且
\[
1,\ y,\ \dots,\ y^{2N-1}
\]
张成整个空间。
对这组函数作 Gram--Schmidt 正交化，得到一组正交归一多项式基。
在该基下，\(M_y\) 表现为一唯一的 \(2N\times 2N\) Jacobi 矩阵 \(J_{\varphi,N}\)。

若 \(U_N:L^2(\mu_{\varphi,N})\to \RR^{2N}\) 把这组正交基送到标准基，则
\[
J_{\varphi,N}=U_NM_yU_N^{-1},
\qquad
v_0^{(N)}:=U_N(1).
\]
于是
\[
u_k^{(N)}
=
\langle 1,M_y^k1\rangle_{L^2(\mu_{\varphi,N})}
=
\langle v_0^{(N)},(J_{\varphi,N})^k v_0^{(N)}\rangle.
\]
由于纯点乘法算子的谱恰为支撑点集，故
\[
\sigma(J_{\varphi,N})=\operatorname{supp}\mu_{\varphi,N},
\]
即为所述两条有限 reciprocal-square ladders。
最后，
\[
m_{\varphi,N}(z)
=
\langle 1,(I-zM_y)^{-1}1\rangle_{L^2(\mu_{\varphi,N})}
=
\int\frac{\dd\mu_{\varphi,N}(y)}{1-yz},
\]
再把 \(\mu_{\varphi,N}\) 展开即可得到有理 Stieltjes 闭式。
证毕。
\end{proof}

\begin{theorem}[黄金双正弦行列式的 finite-rank 局部一致截断]\label{thm:xi-time-part9zbg-foldpi-finite-rank-fredholm-local-uniform}
定义
\[
D_{\varphi,N}(z)
:=
\det(I-zJ_{\varphi,N})
=
\prod_{n=1}^{N}
\left(1-\frac{z}{4\pi^2n^2}\right)
\left(1-\frac{\varphi^2 z}{4\pi^2n^2}\right).
\]
则 \(D_{\varphi,N}\to D_\varphi\) 在 \(\CC\) 的每个紧集上一致收敛。
更精确地，若 \(K\Subset \CC\setminus \mathcal Z(D_\varphi)\) 且
\[
R:=\sup_{z\in K}|z|,
\]
则存在 \(N_0(K)\) 使得对一切 \(N\ge N_0(K)\)，可在 \(K\) 上取一致的对数分支并有
\[
\sup_{z\in K}
\left|
\log D_\varphi(z)-\log D_{\varphi,N}(z)
\right|
\le
\frac{R(1+\varphi^2)}{2\pi^2}\,\frac{1}{N}.
\]
同时，对每个整数 \(r\ge 1\)，有限秩谱 \(\zeta\) 函数满足
\[
\zeta_{J_{\varphi,N}}(r):=\Tr(J_{\varphi,N}^r)
=
(4\pi^2)^{-r}(1+\varphi^{2r})\sum_{n=1}^{N}n^{-2r},
\]
并且
\[
0<
\zeta_{J_\varphi}(r)-\zeta_{J_{\varphi,N}}(r)
\le
(4\pi^2)^{-r}(1+\varphi^{2r})\frac{N^{1-2r}}{2r-1}.
\]
\end{theorem}
\begin{proof}
第一条乘积公式由定理 \ref{thm:xi-time-part9zbg-foldpi-finite-rank-jacobi-fredholm-realization} 的显式谱立即得到。
对紧集 \(K\subset\CC\)，记 \(R:=\sup_{z\in K}|z|\)。
则
\[
\sum_{n>N}\sup_{z\in K}
\left(
\left|\frac{z}{4\pi^2n^2}\right|
+
\left|\frac{\varphi^2 z}{4\pi^2n^2}\right|
\right)
\le
\frac{R(1+\varphi^2)}{4\pi^2}
\sum_{n>N}n^{-2},
\]
而右端随 \(N\to\infty\) 收敛到 \(0\)。
于是尾积
\[
\prod_{n>N}
\left(1-\frac{z}{4\pi^2n^2}\right)
\left(1-\frac{\varphi^2 z}{4\pi^2n^2}\right)
\]
在 \(K\) 上一致收敛到 \(1\)，故 \(D_{\varphi,N}\to D_\varphi\) 局部一致。

若进一步 \(K\Subset\CC\setminus \mathcal Z(D_\varphi)\)，则可取 \(N_0(K)\) 使对 \(N\ge N_0(K)\) 与一切 \(n>N\) 都有
\[
\sup_{z\in K}
\left(
\left|\frac{z}{4\pi^2n^2}\right|
\vee
\left|\frac{\varphi^2 z}{4\pi^2n^2}\right|
\right)
\le
\frac12.
\]
于是利用
\[
|\log(1-w)|\le 2|w|
\qquad (|w|\le 1/2),
\]
便得到
\[
\sup_{z\in K}
\left|
\log\frac{D_\varphi(z)}{D_{\varphi,N}(z)}
\right|
\le
2R\sum_{n>N}
\left(
\frac{1}{4\pi^2n^2}
+
\frac{\varphi^2}{4\pi^2n^2}
\right)
\le
\frac{R(1+\varphi^2)}{2\pi^2}\,\frac1N.
\]

对谱 \(\zeta\) 函数，只需把定理 \ref{thm:xi-time-part9zbg-foldpi-finite-rank-jacobi-fredholm-realization} 的显式谱代入，得到
\[
\zeta_{J_{\varphi,N}}(r)
=
\sum_{n=1}^{N}\left(\frac{1}{4\pi^2n^2}\right)^r
+
\sum_{n=1}^{N}\left(\frac{\varphi^2}{4\pi^2n^2}\right)^r
=
(4\pi^2)^{-r}(1+\varphi^{2r})\sum_{n=1}^{N}n^{-2r}.
\]
而
\[
\zeta_{J_\varphi}(r)-\zeta_{J_{\varphi,N}}(r)
=
(4\pi^2)^{-r}(1+\varphi^{2r})\sum_{n>N}n^{-2r}.
\]
再用积分判别
\[
\sum_{n>N}n^{-2r}
\le
\int_N^\infty x^{-2r}\,\dd x
=
\frac{N^{1-2r}}{2r-1}
\]
即得尾项估计。证毕。
\end{proof}

\begin{theorem}[H.111 Bernoulli 奇层的 finite-rank 单调逼近]\label{thm:xi-time-part9zbg-foldpi-finite-rank-oddlayer-monotone-approx}
\leanverified{paper\_xi\_time\_part9zbg\_foldpi\_finite\_rank\_oddlayer\_monotone\_approx}
对每个 \(r\ge 1\) 与 \(N\ge 1\)，定义 finite-rank odd-layer 系数
\[
\alpha_r^{(N)}
:=
(-1)^{r-1}\,
4(2r-2)!\,
\frac{\varphi^{-1}+\varphi^{2r-3}}{1+\varphi^{2r}}\,
\zeta_{J_{\varphi,N}}(r).
\]
则有显式闭式
\[
\alpha_r^{(N)}
=
(-1)^{r-1}\,
4(2r-2)!\,(4\pi^2)^{-r}
\bigl(\varphi^{-1}+\varphi^{2r-3}\bigr)
\sum_{n=1}^{N}n^{-2r},
\]
并且对每个固定 \(r\) 都有
\[
0<
(-1)^{r-1}\alpha_r^{(1)}
<
(-1)^{r-1}\alpha_r^{(2)}
<
\cdots
<
(-1)^{r-1}\alpha_r,
\]
其中 \(\alpha_r\) 正是 H.111 的 exact odd-layer coefficient。
更精确地，
\[
0<
(-1)^{r-1}\bigl(\alpha_r-\alpha_r^{(N)}\bigr)
\le
\frac{4(2r-2)!}{(4\pi^2)^r}
\frac{\varphi^{-1}+\varphi^{2r-3}}{2r-1}\,
N^{1-2r}.
\]
\end{theorem}
\begin{proof}
将定理 \ref{thm:xi-time-part9zbg-foldpi-finite-rank-fredholm-local-uniform} 的显式公式
\[
\zeta_{J_{\varphi,N}}(r)
=
(4\pi^2)^{-r}(1+\varphi^{2r})\sum_{n=1}^{N}n^{-2r}
\]
代入定义即可得到 \(\alpha_r^{(N)}\) 的闭式。

另一方面，由定理 \ref{thm:xi-time-part9zbe-foldpi-compact-determinant-sampling}，
\[
\alpha_r
=
(-1)^{r-1}\,
4(2r-2)!\,
\frac{\varphi^{-1}+\varphi^{2r-3}}{1+\varphi^{2r}}\,
\zeta_{J_\varphi}(r).
\]
再由定理 \ref{thm:xi-time-part9zbg-foldpi-finite-rank-fredholm-local-uniform} 中的严格尾差
\[
\zeta_{J_\varphi}(r)-\zeta_{J_{\varphi,N}}(r)>0,
\]
立得
\[
0<
(-1)^{r-1}\bigl(\alpha_r-\alpha_r^{(N)}\bigr).
\]
同时利用同一定理中的尾界便得到
\[
0<
(-1)^{r-1}\bigl(\alpha_r-\alpha_r^{(N)}\bigr)
\le
\frac{4(2r-2)!}{(4\pi^2)^r}
\frac{\varphi^{-1}+\varphi^{2r-3}}{2r-1}\,
N^{1-2r}.
\]
故 \(((-1)^{r-1}\alpha_r^{(N)})_{N\ge 1}\) 对固定 \(r\) 严格递增并收敛到 \((-1)^{r-1}\alpha_r\)。
证毕。
\end{proof}

\begin{theorem}[finite-rank 塔中的 Hankel 正性阈值与单调收敛]\label{thm:xi-time-part9zbg-foldpi-finite-rank-hankel-threshold}
对固定 \(n\ge 0\)，定义截断 Hankel 矩阵
\[
H_n^{(N)}:=(u_{i+j}^{(N)})_{0\le i,j\le n},
\qquad
H_n:=(u_{i+j})_{0\le i,j\le n}.
\]
则有
\[
H_n^{(N)}\preceq H_n^{(N+1)}\preceq H_n
\qquad (\forall N\ge 1),
\]
并且只要
\[
2N\ge n+1,
\]
就有 \(H_n^{(N)}\) 严格正定，从而
\[
\det H_n^{(N)}>0.
\]
此外，
\[
\det H_n^{(N)}\uparrow \det H_n>0
\qquad (N\to\infty).
\]
\end{theorem}
\begin{proof}
令
\[
v_n(y):=(1,y,\dots,y^n)^{\mathsf T}.
\]
则
\[
H_n^{(N+1)}-H_n^{(N)}
=
\int v_n(y)v_n(y)^{\mathsf T}\,\dd(\mu_{\varphi,N+1}-\mu_{\varphi,N})(y)\succeq 0,
\]
同理
\[
H_n-H_n^{(N)}
=
\int v_n(y)v_n(y)^{\mathsf T}\,\dd(\mu_\varphi-\mu_{\varphi,N})(y)\succeq 0.
\]
故得到 Loewner 单调性
\[
H_n^{(N)}\preceq H_n^{(N+1)}\preceq H_n.
\]

若 \(2N\ge n+1\)，则 \(\mu_{\varphi,N}\) 至少含有 \(n+1\) 个互异支撑点。
任取 \(c=(c_0,\dots,c_n)\neq 0\)，令
\[
P_c(y):=\sum_{j=0}^{n}c_j y^j.
\]
由于 \(\deg P_c\le n\)，该非零多项式不可能在 \(n+1\) 个互异点上同时为零，故
\[
c^{\mathsf T}H_n^{(N)}c
=
\int P_c(y)^2\,\dd\mu_{\varphi,N}(y)>0.
\]
从而 \(H_n^{(N)}\) 严格正定。

最后，由定理 \ref{thm:xi-time-part9zbg-foldpi-finite-rank-moment-exhaustion}，每个矩条目
\[
u_{i+j}^{(N)}\uparrow u_{i+j},
\]
故
\[
H_n^{(N)}\longrightarrow H_n
\]
逐项收敛。
由行列式的连续性与上面的 Loewner 单调性，得到
\[
\det H_n^{(N)}\uparrow \det H_n.
\]
而 \(\det H_n>0\) 则由定理 \ref{thm:xi-time-part9zbc-foldpi-hankel-strict-positivity} 给出。
证毕。
\end{proof}

\begin{theorem}[finite-depth exactness obstruction]\label{thm:xi-time-part9zbg-foldpi-finite-depth-exactness-obstruction}
\leanverified{paper\_xi\_time\_part9zbg\_foldpi\_finite\_depth\_exactness\_obstruction}
不存在任何 finite-rank scalar spectral host，能够精确承载 H.111 的全体 odd Bernoulli layers。
更具体地，对任意有限 \(N\) 与任意 \(r\ge 1\)，恒有严格不等式
\[
(-1)^{r-1}\alpha_r^{(N)}<(-1)^{r-1}\alpha_r.
\]
因此，exact projective \(\pi\)-tail 不是某个 finite continued fraction 或 finite Jacobi kernel 的改写；其谱宿主必然是无限秩对象。
\end{theorem}
\begin{proof}
前述严格不等式已经由定理 \ref{thm:xi-time-part9zbg-foldpi-finite-rank-oddlayer-monotone-approx} 给出。

更强地，反设存在某个 finite-rank scalar spectral host \(J\) 与循环向量 \(v\)，能够精确承载 H.111 的全部 odd Bernoulli layers。
则通过定理 \ref{thm:xi-time-part9zbe-foldpi-compact-determinant-sampling} 所建立的同一 factorial 采样关系，可得到与 exact tail 对应的矩列仍然必须是
\[
u_k=(-1)^k\frac{\alpha_{k+1}}{(2k)!}
\qquad (k\ge 0).
\]
然而 finite-rank 正自伴 Jacobi 宿主只对应有限支撑测度，因此其矩列的 Hankel 秩有限。
这与定理 \ref{thm:xi-time-part9zbc-foldpi-hankel-strict-positivity} 以及定理 \ref{thm:xi-time-part9zbc-foldpi-zeta-rank-jump} 所给出的“\((u_k)\) 具有无限 Hankel 秩且全部有限 Hankel 主子式严格为正”直接矛盾。
故 exact host 不可能是 finite-rank 对象。

于是，虽然 \((J_{\varphi,N})_{N\ge 1}\) 构成一座从下方单调逼近 exact odd tail 的 finite-depth 塔，但这一塔在任何有限深度都不会稳定为 H.111 的 exact object。
证毕。
\end{proof}

\noindent
因此，H.111 odd Bernoulli tail 在现有审计链下不只是“属于 compact Jacobi/Fredholm 范畴”；更准确地说，它是一个由 finite-depth scalar hosts 严格单调耗尽、却永不在任何有限深度闭合的无限秩对象。每个 \(\mu_{\varphi,N}\) 只给出一个 terminating Stieltjes kernel，每个 \(J_{\varphi,N}\) 只携带有限条 reciprocal-square ladders，每个 \(D_{\varphi,N}\) 只拥有有限零点集，而每个 \(\alpha_r^{(N)}\) 只提供 exact odd tail 的严格下侧逼近。唯有在 \(N\to\infty\) 的严格归纳极限中，exact \(2\pi\)-恢复与其全部 Bernoulli--\(\zeta(2r)\) 谱信息才同时出现。

\endinput
